\begin{theorem}[Reusable Euler-product power identity]\label{mod:finitefield-eulerproduct}
Let $k$ be a finite field and let
\[
 w:\operatorname{MonicPolynomial}(k)\longrightarrow \mathbf C
\]
be a monoid homomorphism.  Write $\mathcal M_n(k)$ for the finite type of
monic polynomials of degree exactly $n$, and $\mathcal P_d(k)$ for its
subtype of irreducible polynomials.  Define
\[
 A_n(w)=\sum_{P\in\mathcal M_n(k)}w(P)
\]
and
\[
 B_r(w)=
 \sum_{d\in\operatorname{divisors}(r)}
 d\sum_{P\in\mathcal P_d(k)}w(P)^{r/d}.
\]
Here $\operatorname{divisors}(0)$ is Mathlib's empty finset, so $B_0(w)=0$.
The corresponding formal power series are
\[
 L_w(T)=\sum_{n\geq 0}A_n(w)T^n,
 \qquad
 C_w(T)=\sum_{n\geq 0}B_{n+1}(w)T^n.
\]
They satisfy the integral formal-power-series identity
\[
 \frac{d}{dT}L_w(T)=L_w(T)C_w(T).
\]

More generally, this derivative identity is proved once for weighted
multisets of arbitrary positive-degree atoms.  If an atom $p$ has degree
$d(p)>0$ and weight $u(p)$, put
\[
 D(s)=\sum_{p\in s}d(p),\qquad W(s)=\prod_{p\in s}u(p).
\]
For every $N$, insertion of a positive block gives an equivalence
\[
 (t,p,h),\quad h>0,\quad D(t)+h d(p)=N
 \quad\longleftrightarrow\quad
 (s,p,j),\quad D(s)=N,\quad j<\operatorname{count}_s(p),
\]
by
\[
 (t,p,h)\longmapsto(t+h[p],p,h-1).
\]
Summing over this equivalence counts every occurrence of every atom and
gives the Newton recurrence
\[
 N\sum_{D(s)=N}W(s)
 =\sum_{i+r=N}
   \left(\sum_{D(t)=i}W(t)\right)
   \left(\sum_{\substack{p,h>0\\h d(p)=r}}
      d(p)u(p)^h\right).
\]
This recurrence is exactly the coefficientwise statement of the displayed
derivative identity.  Its proof takes values in a commutative semiring and
does not divide by a weight or assume that any weight is nonzero.

For the specialization to polynomials, Mathlib's normalized-factor
factorization gives an equivalence
\[
 \operatorname{MonicPolynomial}(k)
 \simeq
 \operatorname{Multiset}\{P:k[X]\mid P\text{ is monic and irreducible}\}.
\]
It preserves both total degree and the product of the weights.  Exact-degree
monic polynomials are finite via Mathlib's equivalences with polynomials of
degree less than $n$ and then with coefficient vectors
$\operatorname{Fin}(n)\to k$.  Finally, for $r>0$, the pairs $(P,h)$ with
$h\deg(P)=r$ are reindexed by $d\mid r$ and
$P\in\mathcal P_d(k)$; this identifies the generic power sum with $B_r(w)$.

Since $A_0(w)=w(1)=1$, if
\[
 A_n(w)=0\qquad(n\geq2),
\]
then coefficient comparison in the derivative identity yields, for every
$r>0$,
\[
 \boxed{-B_r(w)=\bigl(-A_1(w)\bigr)^r.}
\]
This is the exported theorem used unchanged with
$w_{\chi,\psi}$ in the Hasse--Davenport lift and with
$w_{\phi,\psi}$ in the Hasse-function lift.  In particular it allows the
zero value of a multiplicative character on a zero constant term and it
retains repeated closed points through the powers $w(P)^{r/d}$.

\LeanFile{LanglandsFirstMainLemma/FiniteField/EulerProduct.lean}
\lean{LanglandsFirstMainLemma.MonicPolynomialOfDegree}
\lean{LanglandsFirstMainLemma.MonicIrreducibleOfDegree}
\lean{LanglandsFirstMainLemma.monicFactorsEquiv}
\lean{LanglandsFirstMainLemma.GradedEulerProduct.global_weighted_derivative}
\lean{LanglandsFirstMainLemma.MonicWeight.degreeSum}
\lean{LanglandsFirstMainLemma.MonicWeight.closedPointPowerSum}
\lean{LanglandsFirstMainLemma.MonicWeight.closedPointSeries_derivative}
\lean{LanglandsFirstMainLemma.closedPoint_power_identity}
\leanok
\SourceText{Proofs of Theorems thm:HD-lift and thm:Hasse-lift.}
\uses{mod:finitefield-polynomialweights}
\FormalizationContract
\end{theorem}
